\documentclass[11pt]{article}

% -------------------------
% Packages
% -------------------------
\usepackage[a4paper,margin=1in]{geometry}
\usepackage{amsmath, amssymb, amsthm}
\usepackage{mathtools}
\usepackage{graphicx}
\usepackage[
  colorlinks=true,
  linkcolor=black,
  urlcolor=black,
  citecolor=black
]{hyperref}

\usepackage{xcolor}
\usepackage{enumitem}
\usepackage{booktabs}
% -------------------------
% Section Spacing + Divider
% -------------------------
\usepackage{titlesec}
\usepackage{tikz}

% Section spacing (clear separation)
\titlespacing*{\section}
{0pt}                % left
{5.0ex plus 1.5ex}   % space BEFORE section title
{2.8ex plus 0.5ex}   % space AFTER section title


% Subsection spacing
\titlespacing*{\section}
{0pt}{2.8ex plus .6ex}{0.8ex}


% Subsubsection spacing (compact)
\titlespacing*{\subsubsection}
{0pt}{2.0ex plus .5ex minus .2ex}{1.0ex plus .2ex}

% Dashed divider after each section
\newcommand{\sectionline}{
  \vspace{0.3em}
  \noindent\begin{tikzpicture}
    \draw[dashed, gray] (0,0) -- (\linewidth,0);
  \end{tikzpicture}
  \vspace{0.4em}
}


% Automatically insert divider after every section
\titleformat{\section}
  {\normalfont\Large\bfseries}
  {\thesection}{1em}{}
  [\sectionline]


% -------------------------
% Code Formatting
% -------------------------
\usepackage{listings}
\lstset{
    language=C++,
    basicstyle=\ttfamily\small,
    keywordstyle=\color{blue},
    stringstyle=\color{red},
    commentstyle=\color{gray},
    numbers=left,
    numberstyle=\tiny,
    stepnumber=1,
    numbersep=5pt,
    frame=single,
    breaklines=true,
    tabsize=2,
    showstringspaces=false
}

% -------------------------
% Title
% -------------------------
\title{\textbf{Competitive Programming Maths}}
\author{}
\date{}
\setcounter{tocdepth}{2}


\begin{document}

\maketitle
\tableofcontents
\newpage

% =====================================================
% PASTE ALL CONTENT BELOW THIS LINE
% =====================================================

\section{Number Bases and Binary Mathematics}

\subsection{What Are Number Bases?}

\subsubsection{Definition}
A \textbf{number base} (or radix) is the number of distinct digits used to represent numbers, including zero.
In base $b$, each digit position represents a power of $b$.

\[
(d_k d_{k-1} \ldots d_1 d_0)_b = \sum_{i=0}^{k} d_i \cdot b^i
\]

\subsection{Binary Numbers}

\subsubsection{Representation}
Binary numbers use base 2.

\[
(101101)_2 = 1\cdot2^5 + 0\cdot2^4 + 1\cdot2^3 + 1\cdot2^2 + 0\cdot2^1 + 1\cdot2^0 = 45
\]

\subsubsection{Key Properties}
\begin{itemize}
    \item Each bit represents a power of 2
    \item $k$ bits represent values from $0$ to $2^k - 1$
    \item LSB decides parity
\end{itemize}

\subsection{Bitwise Operators}

\begin{center}
\begin{tabular}{ccc}
\toprule
Operator & Symbol & Meaning \\
\midrule
AND & \texttt{\&} & Both bits must be 1 \\
OR & \texttt{|} & Any bit is 1 \\
XOR & \texttt{\^} & Bits differ \\
NOT & \texttt{\~} & Flip bits \\
Left Shift & \texttt{<<} & Multiply by $2^k$ \\
Right Shift & \texttt{>>} & Divide by $2^k$ \\
\bottomrule
\end{tabular}
\end{center}

\subsection{Binary Tricks}

\subsubsection{Power of Two Check}
\[
n > 0 \land (n \& (n-1)) = 0
\]

\textbf{Why it works:}  
A power of two has exactly one set bit. Subtracting 1 flips that bit and all bits after it, so AND becomes zero.

---

\subsubsection{Lowest Set Bit}
\[
n \& (-n)
\]

\textbf{Meaning:}  
Extracts the value of the least significant set bit.

\textbf{Example:}
\[
n = 12 = (1100)_2 \Rightarrow n \& (-n) = 4
\]

---

\subsubsection{Important XOR Identities}

\begin{itemize}
    \item $a \oplus a = 0$
    \item $a \oplus 0 = a$
    \item $a \oplus b \oplus a = b$
    \item XOR is \textbf{commutative} and \textbf{associative}
\end{itemize}

---

\subsubsection{AND Operator Properties}

\begin{itemize}
    \item $a \& a = a$
    \item $a \& 0 = 0$
    \item $a \& b \le \min(a, b)$
\end{itemize}

\textbf{Observation:}  
AND only keeps bits that are set in \textbf{both} numbers.

---

\subsubsection{OR Operator Properties}

\begin{itemize}
    \item $a | a = a$
    \item $a | 0 = a$
    \item $a | b \ge \max(a, b)$
\end{itemize}

\textbf{Observation:}  
OR sets a bit if it appears in \textbf{either} number.

---

\subsubsection{Bitwise NOT}

\[
\sim a = -a - 1
\]

\textbf{Why this holds:}  
Integers are stored in two's complement.

\textbf{Example:}
\[
a = 5 \Rightarrow \sim a = -6
\]

---

\subsubsection{Useful CP Identities}

\begin{itemize}
    \item $(a \oplus b) \oplus c = a \oplus (b \oplus c)$
    \item $a + b = (a \oplus b) + 2 \cdot (a \& b)$
    \item Removing lowest set bit:
    \[
    n = n \& (n-1)
    \]
\end{itemize}

---


\subsection{C++ Built-in Functions}

\subsubsection{Popcount}
\begin{lstlisting}
int x = 13;
int bits = __builtin_popcount(x); // 3
\end{lstlisting}

\subsubsection{Leading / Trailing Zeros}
\begin{lstlisting}
int x = 16;
__builtin_ctz(x); // 4
__builtin_clz(x); // depends on bit-width
\end{lstlisting}

\subsection{What is $\log_2$?}

\subsubsection{Definition}
\[
\log_2(n) = \text{the power to which 2 must be raised to obtain } n
\]

\subsubsection{CP Interpretation}
\begin{itemize}
    \item Number of bits required to represent $n$
    \item Number of times we can divide $n$ by 2 before it becomes zero
\end{itemize}

\subsubsection{Properties of $\log_2$}

For positive numbers $a, b$:

\begin{itemize}
    \item $\log_2(1) = 0$
    \item $\log_2(2^k) = k$
    \item $\log_2(ab) = \log_2(a) + \log_2(b)$
    \item $\log_2\!\left(\frac{a}{b}\right) = \log_2(a) - \log_2(b)$
    \item $\log_2(a^k) = k \log_2(a)$
\end{itemize}

\subsubsection{Important Inequality}
\[
2^{\lfloor \log_2(n) \rfloor} \le n < 2^{\lfloor \log_2(n) \rfloor + 1}
\]

\subsubsection{Relation with Binary Length}
The number of bits required to represent a positive integer $n$ is:
\[
\lfloor \log_2(n) \rfloor + 1
\]

\subsection{Edge Cases}
\begin{itemize}
    \item \texttt{\_\_builtin\_clz(0)} is undefined
    \item Shifting negative integers is implementation-defined
    \item Avoid floating-point \texttt{log2} due to precision issues
    \item Prefer bitwise operations for exact results
\end{itemize}


% =====================================================
\section{Data Types and Precision in C++}
% =====================================================

\subsection{Why Data Types Matter}
Choosing the correct data type is critical in competitive programming because:
\begin{itemize}
    \item Overflow bugs silently produce wrong answers
    \item Bitwise operations depend on exact binary representation
    \item Floating-point precision errors can break geometry and math problems
\end{itemize}

% =====================================================
\subsection{Integer Data Types}
% =====================================================

\subsubsection{Signed Integers}

\begin{center}
\begin{tabular}{cccc}
\toprule
Type & Bits & Decimal Range & Binary Representation \\
\midrule
\texttt{int} & 32 &
$[-2^{31}, 2^{31}-1]$ &
Two's Complement \\
\texttt{long long} & 64 &
$[-2^{63}, 2^{63}-1]$ &
Two's Complement \\
\bottomrule
\end{tabular}
\end{center}

\paragraph{Two's Complement Representation}
Signed integers are stored using \textbf{two's complement}:
\begin{itemize}
    \item MSB indicates sign
    \item Negative numbers are represented as:
    \[
    -x = \sim x + 1
    \]
\end{itemize}

\subsubsection{Unsigned Integers}

\begin{center}
\begin{tabular}{cccc}
\toprule
Type & Bits & Decimal Range & Binary Range \\
\midrule
\texttt{unsigned int} & 32 &
$[0, 2^{32}-1]$ &
$[0, 111\ldots111_2]$ \\
\texttt{unsigned long long} & 64 &
$[0, 2^{64}-1]$ &
$[0, 111\ldots111_2]$ \\
\bottomrule
\end{tabular}
\end{center}

\paragraph{CP Rule}
Use \texttt{long long} by default.  

% =====================================================
\subsection{Floating Point Data Types}
% =====================================================

\subsubsection{Types and Precision}

\begin{center}
\begin{tabular}{cccc}
\toprule
Type & Size & Significant Digits & Exponent Range \\
\midrule
\texttt{float} & 32-bit & $\sim 6$ & $\sim 10^{\pm38}$ \\
\texttt{double} & 64-bit & $\sim 15$ & $\sim 10^{\pm308}$ \\
\texttt{long double} & 80-bit (gcc) & $\sim 18$ & Larger \\
\bottomrule
\end{tabular}
\end{center}

\subsubsection{IEEE 754 Representation (Conceptual)}
A floating-point number is stored as:
\[
(-1)^{sign} \times mantissa \times 2^{exponent}
\]

Because mantissa has limited bits, many decimals cannot be represented exactly.

% =====================================================
\subsection{Floating Point Inaccuracies (With Examples)}
% =====================================================

\subsubsection{Why Floating Point Inaccuracies Occur}
Many decimal numbers cannot be represented exactly in binary.

\[
0.1_{10} = 0.00011001100110011\ldots_2
\]

Since the binary expansion is infinite, it must be truncated, leading to approximation errors.

\subsubsection{Example 1: Unexpected Output}

\begin{lstlisting}
double x = 0.1 + 0.2;
cout << fixed << setprecision(17) << x;
\end{lstlisting}

\textbf{Output:}
\[
0.30000000000000004
\]

\subsubsection{Example 2: Equality Comparison Failure}

\begin{lstlisting}
double a = 0.3;
double b = 0.1 + 0.2;

if (a == b)
    cout << "Equal";
else
    cout << "Not Equal";
\end{lstlisting}

\textbf{Output:}
\[
\texttt{Not Equal}
\]

% =====================================================
\subsection{Comparing Floating Point Numbers}
% =====================================================

\subsubsection{Why \texttt{==} is Dangerous}
\begin{itemize}
    \item Floating-point numbers are approximations
    \item Exact equality rarely holds
\end{itemize}

\subsubsection{Epsilon-Based Comparison}

\paragraph{Definition}
\[
\varepsilon = \text{a very small positive number}
\]

\paragraph{Rule}
\[
|a - b| < \varepsilon \Rightarrow a \approx b
\]

\paragraph{Typical Values}
\begin{itemize}
    \item \texttt{1e-9} for \texttt{double}
    \item \texttt{1e-12} or smaller for \texttt{long double}
\end{itemize}

\subsubsection{C++ Implementation}
\begin{lstlisting}
const double EPS = 1e-9;

bool equal(double a, double b) {
    return fabs(a - b) < EPS;
}
\end{lstlisting}

% =====================================================
\subsection{Competitive Programming Notes}
% =====================================================

\begin{itemize}
    \item Use \texttt{long long} everywhere by default to avoid overflow
    \item \texttt{int} operations are faster than \texttt{long long}, but safety matters more
    \item Prefer integer arithmetic over floating-point whenever possible
    \item Comparing floating-point numbers using \texttt{==} leads to incorrect results
    \item \texttt{\_\_builtin\_clz(0)} and similar functions are undefined
    \item Shifting negative integers is implementation-defined
\end{itemize}

% =====================================================
\subsection{Why \texttt{int} can give Runtime Error}
% =====================================================

\subsubsection{Problem Reference (CP Context)}

This issue appears in the following real competitive programming problem:

\begin{itemize}
    \item \textbf{Divide Two Integers}
    \item Platform: LeetCode
    \item Link: \url{https://leetcode.com/problems/divide-two-integers/description/}
\end{itemize}

\subsubsection{The Problematic Case}

Consider the division:
\[
-2^{31} \div (-1)
\]

\textbf{Mathematically:}
\[
-2^{31} \div (-1) = 2^{31}
\]

However,
\[
2^{31} = 2147483648 > \texttt{INT\_MAX}
\]

\subsubsection{Why This Causes a Runtime Error}

A 32-bit signed integer (\texttt{int}) has the range:
\[
[-2^{31},\ 2^{31}-1]
\]

The value $2147483648$ cannot be represented using \texttt{int}.  
In C++, this results in \textbf{undefined behavior}.

\subsubsection{Failing Code Example}

\begin{lstlisting}
int divide(int dividend, int divisor) {
    return dividend / divisor;
}
\end{lstlisting}

\textbf{Input:}
\[
\texttt{dividend = -2147483648,\ divisor = -1}
\]

\textbf{Observed Result:}
\begin{itemize}
    \item Runtime Error
    \item Integer overflow
    \item Undefined behavior
\end{itemize}

\subsubsection{Correct Approach Using \texttt{long long}}

\begin{lstlisting}
int divide(int dividend, int divisor) {
    return min(2147483647LL,(long long)dividend / divisor));
}
\end{lstlisting}

% =====================================================
\section{Modular Arithmetic and Binary Exponentiation}
% =====================================================

\subsection{What is Modular Arithmetic?}

\subsubsection{Definition}
For a positive integer $m$, we say:
\[
a \equiv b \pmod{m}
\]
if $(a - b)$ is divisible by $m$.

In competitive programming, answers are often required modulo a fixed number, most commonly:
\[
m = 10^9 + 7 \quad \text{or} \quad m = 998244353
\]


\subsubsection{Why Modulo is Used}
\begin{itemize}
    \item Prevents integer overflow
    \item Keeps values within fixed bounds
    \item Required explicitly in many problems
\end{itemize}

% =====================================================
\subsection{Basic Properties of Modulo}
% =====================================================

For any integers $a, b$:

\begin{align*}
(a + b) \bmod m &= ((a \bmod m) + (b \bmod m)) \bmod m \\
(a - b) \bmod m &= ((a \bmod m) - (b \bmod m) + m) \bmod m \\
(a \cdot b) \bmod m &= ((a \bmod m) \cdot (b \bmod m)) \bmod m
\end{align*}

\subsubsection{Important Warning}
\[
(a / b) \bmod m \neq (a \bmod m) / (b \bmod m)
\]

Division under modulo requires a \textbf{modular inverse}.

% =====================================================
\subsection{Binary Exponentiation}
% =====================================================

\subsubsection{Problem}
Compute:
\[
a^b \bmod m
\]
efficiently.

A naive approach takes $O(b)$ time, which is too slow.

\subsubsection{Idea}
Represent the exponent in binary and repeatedly square the base.

\subsubsection{Algorithm}
\begin{lstlisting}
long long binpow(long long a, long long b, long long m) {
    long long res = 1;
    a %= m;
    while (b) {
        if (b & 1) res = (res * a) % m;
        a = (a * a) % m;
        b >>= 1;
    }
    return res;
}
\end{lstlisting}

\subsubsection{Time Complexity}
\[
O(\log b)
\]

% =====================================================
\subsection{Modular Inverse}
% =====================================================

\subsubsection{Definition}
The modular inverse of $a$ modulo $m$ is a number $x$ such that:
\[
a \cdot x \equiv 1 \pmod{m}
\]

\subsubsection{Existence Condition}
A modular inverse exists if and only if:
\[
\gcd(a, m) = 1
\]

% =====================================================
\subsection{Modular Inverse Using Fermat's Little Theorem}
% =====================================================

\subsubsection{Theorem}
If $m$ is prime and $\gcd(a, m) = 1$, then:
\[
a^{m-1} \equiv 1 \pmod{m}
\]

\subsubsection{Proof}
From Fermat's Little Theorem:
\[
a^{m-1} \equiv 1 \pmod{m}
\]

Multiplying both sides by $a^{-1}$:
\[
a^{m-2} \equiv a^{-1} \pmod{m}
\]

Hence,
\[
\boxed{a^{-1} \equiv a^{m-2} \pmod{m}}
\]

% =====================================================
\subsection{Reference Implementation}
% =====================================================

\subsubsection{Modular Power}

\begin{lstlisting}
long long modpow(long long a, long long b, long long m) {
    long long res = 1;
    a %= m;
    while (b) {
        if (b & 1) res = (res * a) % m;
        a = (a * a) % m;
        b >>= 1;
    }
    return res;
}
\end{lstlisting}

\subsubsection{Modular Inverse}

\begin{lstlisting}
long long inv(long long a, long long m) {
    return modpow(a, m - 2, m);
}
\end{lstlisting}

\subsubsection{Division Under Modulo}

\[
\frac{a}{b} \bmod m = a \cdot b^{-1} \bmod m
\]

\begin{lstlisting}
long long divide(long long a, long long b, long long m) {
    return (a % m) * inv(b, m) % m;
}
\end{lstlisting}


% =====================================================
\subsection{Common Pitfalls}
% =====================================================

\begin{itemize}
    \item Modular inverse does not exist if $\gcd(a, m) \neq 1$
    \item Fermat's method works only for prime moduli
    \item Forgetting modulo during multiplication causes overflow
    \item Direct division under modulo is incorrect
\end{itemize}

% =====================================================
\section{Greatest Common Divisor (GCD) and Euclidean Algorithm}
% =====================================================

\subsection{What is GCD?}

\subsubsection{Definition}
The \textbf{Greatest Common Divisor (GCD)} of two integers $a$ and $b$,
denoted as $\gcd(a, b)$, is the largest positive integer that divides both
$a$ and $b$.

\[
\gcd(a, b) = \max \{ d \mid d \mid a \text{ and } d \mid b \}
\]

\subsubsection{Examples}
\[
\gcd(12, 18) = 6
\]
\[
\gcd(7, 13) = 1
\]

If $\gcd(a, b) = 1$, the numbers are called \textbf{coprime}.

% =====================================================
\subsection{Properties of GCD}
% =====================================================

\begin{itemize}
    \item $\gcd(a, 0) = |a|$
    \item $\gcd(0, 0)$ is undefined
    \item $\gcd(a, b) = \gcd(b, a)$
    \item $\gcd(a, b) = \gcd(b, a \bmod b)$
    \item $\gcd(a, b) \mid a$ and $\gcd(a, b) \mid b$
    \item If $d \mid a$ and $d \mid b$, then $d \mid \gcd(a, b)$
\end{itemize}

\subsubsection{Relation with LCM}
\[
\gcd(a, b) \cdot \text{lcm}(a, b) = |a \cdot b|
\]

% =====================================================
\subsection{Euclidean Algorithm}
% =====================================================

\subsubsection{Idea}
The Euclidean Algorithm is based on the observation:
\[
\gcd(a, b) = \gcd(b, a \bmod b)
\]

We repeatedly replace the larger number by the remainder until it becomes zero.

\subsubsection{Algorithm Steps}
\begin{enumerate}
    \item If $b = 0$, return $a$
    \item Otherwise, compute $\gcd(b, a \bmod b)$
\end{enumerate}

% =====================================================
\subsection{Proof of Euclidean Algorithm}
% =====================================================

Let $a = bq + r$, where $r = a \bmod b$.

\begin{itemize}
    \item Any number dividing both $a$ and $b$ also divides $r = a - bq$
    \item Any number dividing both $b$ and $r$ also divides $a$
\end{itemize}

Hence, the set of common divisors of $(a, b)$ and $(b, r)$ is the same.

\[
\Rightarrow \gcd(a, b) = \gcd(b, a \bmod b)
\]

This justifies the correctness of the algorithm.

% =====================================================
\subsection{Time Complexity}
% =====================================================

Each step reduces the second argument.

\[
\text{Time Complexity} = O(\log \min(a, b))
\]

This makes the Euclidean Algorithm extremely efficient even for large numbers.

% =====================================================
\subsection{C++ Implementation}
% =====================================================

\subsubsection{Recursive Version}

\begin{lstlisting}
long long gcd(long long a, long long b) {
    if (b == 0)
        return a;
    return gcd(b, a % b);
}
\end{lstlisting}

\subsubsection{Iterative Version}

\begin{lstlisting}
long long gcd(long long a, long long b) {
    while (b) {
        a %= b;
        swap(a, b);
    }
    return a;
}
\end{lstlisting}

\subsubsection{Using C++ Built-in Function}

\begin{lstlisting}
long long g = __gcd(a, b); // or std::gcd(a, b) in C++17+
\end{lstlisting}

% =====================================================
\subsection{Practice Problem: Yet Another Array Problem}
% =====================================================

\subsubsection{Problem Reference}

\begin{itemize}
    \item \textbf{Problem}: D. Yet Another Array Problem
    \item Platform: Codeforces
    \item Link: \url{https://codeforces.com/contest/2167/problem/D}
\end{itemize}

\subsubsection{Problem Summary}

You are given an array $a$ of length $n$.
Find the smallest integer $x \ (2 \le x \le 10^{18})$ such that:
\[
\exists i \text{ with } \gcd(a_i, x) = 1
\]

If no such $x$ exists, output $-1$.

---

\subsubsection{Key Observation}

If $\gcd(a_i, x) = 1$, then:
\begin{itemize}
    \item $x$ shares \textbf{no common prime factor} with $a_i$
\end{itemize}

So instead of reasoning about $x$ directly, we reason about its \textbf{prime divisors}.

---

\subsubsection{Important Claim}

Let $x$ be a valid answer and let $p$ be the smallest prime divisor of $x$.

Then:
\[
\gcd(p, a_i) = 1 \quad \text{for some } i
\]

\subsubsection{Proof of the Claim}

Let:
\[
x = p_1^{e_1} \cdot p_2^{e_2} \cdots p_k^{e_k}
\]

and:
\[
a_i = q_1^{f_1} \cdot q_2^{f_2} \cdots
\]

We know:
\[
\gcd(x, a_i) = \prod p_j^{\min(e_j, f_j)} = 1
\]

Hence, for every prime divisor $p_j$ of $x$, we must have:
\[
f_j = 0
\]

In particular, for the smallest prime divisor $p$ of $x$:
\[
\gcd(p, a_i) = 1
\]

Thus, it is sufficient to check only \textbf{prime numbers} as candidates for $x$.

---

\subsubsection{Why Brute Force Over Small Primes Works}

We are asked for the \textbf{smallest} such $x$.

Observe:
\[
2 \cdot 3 \cdot 5 \cdot 7 \cdots 53 > 10^{18}
\]

So any number $x \le 10^{18}$ must have a prime factor $\le 53$.

Therefore:
\begin{itemize}
    \item It is sufficient to check primes up to $53$
    \item If none work, answer is $-1$
\end{itemize}

---

\subsubsection{Algorithm}

\begin{enumerate}
    \item Generate all primes $\le 53$
    \item For each prime $p$:
    \begin{itemize}
        \item Check if $\exists a_i$ such that $\gcd(a_i, p) = 1$
        \item If yes, output $p$
    \end{itemize}
    \item If no prime satisfies the condition, output $-1$
\end{enumerate}

---

\subsubsection{Time Complexity}

\[
O(\text{number of primes} \times n) \approx O(16 \cdot n)
\]

This easily fits within constraints.

---

% =====================================================
\section{Prime Numbers and Factorization Techniques}
% =====================================================

\subsection{Prime Numbers}

\subsubsection{Definition}
A \textbf{prime number} is an integer greater than $1$ that has exactly two
positive divisors:
\[
1 \text{ and itself}
\]

Numbers greater than $1$ that are not prime are called \textbf{composite}.


% =====================================================
\subsection{Primality Testing}
% =====================================================

\subsubsection{Brute Force Primality Test}

\paragraph{Idea}
Check whether $n$ is divisible by every integer from $2$ to $n-1$.

\paragraph{Why It Works}
If $n$ has a divisor other than $1$ and itself, it must appear in this range.

\paragraph{Time Complexity}
\[
O(n)
\]

\paragraph{Code}
\begin{lstlisting}
bool isPrime(long long n) {
    if (n <= 1) return false;
    for (long long i = 2; i < n; i++)
        if (n % i == 0)
            return false;
    return true;
}
\end{lstlisting}

---

\subsubsection{Optimization 1: Checking up to $\sqrt{n}$}

\paragraph{Key Observation}
If $n = a \cdot b$, then at least one of $a, b \le \sqrt{n}$.

\paragraph{What This Improves}
Reduces checks from $n$ to $\sqrt{n}$.

\paragraph{Time Complexity}
\[
O(\sqrt{n})
\]

\paragraph{Code}
\begin{lstlisting}
bool isPrime(long long n) {
    if (n <= 1) return false;
    for (long long i = 2; i * i <= n; i++)
        if (n % i == 0)
            return false;
    return true;
}
\end{lstlisting}

---

\subsubsection{Optimization 2: $6k \pm 1$ Method}

\paragraph{Observation}
All integers can be written as:
\[
6k,\; 6k \pm 1,\; 6k \pm 2,\; 6k + 3
\]

Numbers divisible by $2$ or $3$ are composite.
Thus, all primes greater than $3$ are of the form $6k \pm 1$.

\paragraph{Time Complexity}
\[
O(\sqrt{n}/6)
\]

\paragraph{Code}
\begin{lstlisting}
bool isPrime(long long n) {
    if (n <= 1) return false;
    if (n <= 3) return true;
    if (n % 2 == 0 || n % 3 == 0) return false;

    for (long long i = 5; i * i <= n; i += 6)
        if (n % i == 0 || n % (i + 2) == 0)
            return false;

    return true;
}
\end{lstlisting}

% =====================================================
\subsection{Generating All Prime Numbers up to a Limit}
% =====================================================

\subsubsection{Brute Force Method}

\paragraph{Idea}
Check primality of every number from $2$ to \texttt{limit}.

\paragraph{Why It Is Slow}
Each primality check costs $O(\sqrt{n})$.

\paragraph{Time Complexity}
\[
O(n\sqrt{n})
\]

---

\subsubsection{Optimized Method: Sieve of Eratosthenes}

\paragraph{Core Idea}
If $i$ is prime, all multiples of $i$ are composite.

\paragraph{Why Start from $i^2$}
Smaller multiples already have smaller prime factors.

\paragraph{What It Gives}
All prime numbers $\le$ \texttt{limit}.

\paragraph{Time Complexity}
\[
O(n \log \log n)
\]

\paragraph{Code}
\begin{lstlisting}
vector<long long> generatePrimesUpTo(long long limit) {
    vector<bool> isPrime(limit + 1, true);
    vector<long long> primes;

    isPrime[0] = isPrime[1] = false;
    for (long long i = 2; i * i <= limit; i++) {
        if (isPrime[i]) {
            for (long long j = i * i; j <= limit; j += i)
                isPrime[j] = false;
        }
    }

    for (long long i = 2; i <= limit; i++)
        if (isPrime[i]) primes.push_back(i);

    return primes;
}
\end{lstlisting}

% =====================================================
\subsection{Prime Factorization of a Single Number}
% =====================================================

\subsubsection{Brute Force Factorization}

\paragraph{Idea}
Try dividing $n$ by all numbers from $2$ to $n$.

\paragraph{Time Complexity}
\[
O(n)
\]

---

\subsubsection{Optimized Prime Factorization}

\paragraph{Key Observation}
After removing all factors $\le \sqrt{n}$, at most one prime factor remains.

\paragraph{Time Complexity}
\[
O(\sqrt{n})
\]

\paragraph{Code}
\begin{lstlisting}
vector<long long> primeFactorize(long long n) {
    vector<long long> factors;
    for (long long i = 2; i * i <= n; i++) {
        while (n % i == 0) {
            factors.push_back(i);
            n /= i;
        }
    }
    if (n > 1)
        factors.push_back(n);
    return factors;
}
\end{lstlisting}

% =====================================================
\subsection{Finding All Divisors of a Number}
% =====================================================

\subsubsection{Brute Force Method}

\paragraph{Idea}
Check all numbers from $1$ to $n$.

\paragraph{Time Complexity}
\[
O(n)
\]

---

\subsubsection{Optimized Method}

\paragraph{Key Observation}
Divisors come in pairs:
\[
d \cdot \frac{n}{d} = n
\]

\paragraph{Time Complexity}
\[
O(\sqrt{n})
\]

\paragraph{Code}
\begin{lstlisting}
vector<long long> getAllDivisors(long long n) {
    vector<long long> divs;
    for (long long i = 1; i * i <= n; i++) {
        if (n % i == 0) {
            divs.push_back(i);
            if (i != n / i)
                divs.push_back(n / i);
        }
    }
    sort(divs.begin(), divs.end());
    return divs;
}
\end{lstlisting}

% =====================================================
\subsection{Finding All Divisors of All Numbers up to $N$}
% =====================================================

\subsubsection{Brute Force Idea}

\paragraph{Time Complexity}
\[
O(N^2)
\]

---

\subsubsection{Optimized Divisor Sieve}

\paragraph{Key Idea}
If $i$ divides $j$, then $i$ is a divisor of $j$.

\paragraph{Time Complexity}
\[
O(N \log N)
\]

\paragraph{Code}
\begin{lstlisting}
vector<vector<long long>> allDivisorsUpToN(long long n) {
    vector<vector<long long>> divs(n + 1);
    for (long long i = 1; i <= n; i++)
        for (long long j = i; j <= n; j += i)
            divs[j].push_back(i);
    return divs;
}
\end{lstlisting}

% =====================================================
\subsection{Prime Factors of All Numbers up to $N$}
% =====================================================

\subsubsection{Idea}
If $p$ is prime, then $p$ is a prime factor of all multiples of $p$.

\paragraph{What It Gives}
Prime factor list for every number.

\paragraph{Time Complexity}
\[
O(N \log \log N)
\]

\paragraph{Code}
\begin{lstlisting}
vector<vector<long long>>
primeFactorsUpToN(long long n, const vector<long long>& primes) {
    vector<vector<long long>> factors(n + 1);
    for (auto p : primes)
        for (long long j = p; j <= n; j += p)
            factors[j].push_back(p);
    return factors;
}
\end{lstlisting}

% =====================================================
\section{Smallest Prime Factor (SPF)}
% =====================================================

\subsection{What is the Smallest Prime Factor?}

\subsubsection{Definition}
For every integer $x \ge 2$, the \textbf{Smallest Prime Factor (SPF)} of $x$,
denoted as $\text{spf}[x]$, is the smallest prime number that divides $x$.

Examples:
\[
\text{spf}[2] = 2,\quad
\text{spf}[6] = 2,\quad
\text{spf}[15] = 3,\quad
\text{spf}[17] = 17
\]

If $\text{spf}[x] = x$, then $x$ is prime.

---

\subsection{Why Do We Need SPF?}

\subsubsection{The Problem with Earlier Methods}
\begin{itemize}
    \item Trial division takes $O(\sqrt{n})$ per number
    \item Factoring many numbers becomes too slow
    \item Repeating work across queries is inefficient
\end{itemize}

SPF solves this by \textbf{preprocessing once} and answering queries fast.

---

\subsection{Core Idea Behind SPF}

\subsubsection{Connection with the Sieve of Eratosthenes}
SPF is an extension of the sieve.

\begin{itemize}
    \item Initially, assume $\text{spf}[i] = i$
    \item When a prime $i$ is found, mark it as the smallest divisor
    \item For each multiple $j = i^2, i^2+i, \dots$:
    \begin{itemize}
        \item If $\text{spf}[j]$ is unset, set $\text{spf}[j] = i$
    \end{itemize}
\end{itemize}

This ensures every number gets its smallest prime divisor.

---

\subsection{Why Starting from $i^2$ Works}

If $j < i^2$, then:
\[
j = i \cdot k \quad \Rightarrow \quad k < i
\]

So $j$ would already have a smaller prime factor and its SPF would have been set earlier.

---

\subsection{What Does SPF Give Us?}

Using SPF, we can:
\begin{itemize}
    \item Check if a number is prime in $O(1)$
    \item Factorize a number in $O(\log n)$
    \item Generate all primes up to $N$
    \item Precompute prime factors for all numbers
\end{itemize}

---

\subsection{Time and Space Complexity}

\begin{itemize}
    \item Precomputation: $O(N \log \log N)$
    \item Space: $O(N)$
    \item Primality check: $O(1)$
    \item Factorization of one number: $O(\log n)$
\end{itemize}

---

\subsection{SPF Construction}

\subsubsection{Code}
\begin{lstlisting}
vector<long long> computeSPF(long long n) {
    vector<long long> spf(n + 1);
    for (long long i = 1; i <= n; i++)
        spf[i] = i;

    for (long long i = 2; i * i <= n; i++) {
        if (spf[i] == i) { // i is prime
            for (long long j = i * i; j <= n; j += i) {
                if (spf[j] == j)
                    spf[j] = i;
            }
        }
    }
    return spf;
}
\end{lstlisting}

---

\subsection{Applications of SPF}

\subsubsection{Primality Check}

\begin{lstlisting}
bool isPrime(long long x, const vector<long long>& spf) {
    return x >= 2 && spf[x] == x;
}
\end{lstlisting}

---

\subsubsection{Generating All Primes up to $N$}

\begin{lstlisting}
vector<long long> generatePrimes(long long n,
                                 const vector<long long>& spf) {
    vector<long long> primes;
    for (long long i = 2; i <= n; i++) {
        if (spf[i] == i)
            primes.push_back(i);
    }
    return primes;
}
\end{lstlisting}

---

\subsubsection{Prime Factorization Using SPF}

\begin{lstlisting}
vector<long long> primeFactors(long long n,
                               const vector<long long>& spf) {
    vector<long long> factors;
    while (n > 1) {
        factors.push_back(spf[n]);
        n /= spf[n];
    }
    return factors;
}
\end{lstlisting}

---
% =====================================================
\subsection{Practice Problem: Reducing Fractions}
% =====================================================

\paragraph{Problem Link}
\url{https://codeforces.com/contest/222/problem/C}

\paragraph{Problem Summary}

You are given two arrays representing the numerator and denominator of a fraction.
Each array element contributes multiplicatively.

Your task is to \textbf{reduce the fraction} by cancelling common prime factors
between the numerator and denominator, and output the reduced form.

The constraints allow values up to $10^7$ and many numbers must be factorized.

---

\subsubsection{Core Observation}

Reduction of a fraction depends only on \textbf{prime factor counts}.

For each prime $p$:
\begin{itemize}
    \item Let $a_p$ = total exponent of $p$ in the numerator
    \item Let $b_p$ = total exponent of $p$ in the denominator
\end{itemize}

After reduction:
\[
\text{new } a_p = a_p - \min(a_p, b_p),
\quad
\text{new } b_p = b_p - \min(a_p, b_p)
\]

So the problem becomes:
\begin{itemize}
    \item Efficiently factorize many numbers
    \item Maintain prime exponent counts
    \item Reconstruct reduced numbers
\end{itemize}

---

% =====================================================
\subsubsection{Approach 1: Using Smallest Prime Factor (Recommended)}
% =====================================================

\subsubsubsection{Idea}

\begin{itemize}
    \item Precompute the \textbf{Smallest Prime Factor (SPF)} for all numbers up to $10^7$
    \item Factorize each number in $O(\log n)$ using SPF
    \item Accumulate prime exponent counts for numerator and denominator
    \item Cancel common prime powers
    \item Rebuild the reduced arrays
\end{itemize}

\subsubsubsection{Why This Works Well}

\begin{itemize}
    \item SPF allows factorization by repeatedly dividing by the smallest prime
    \item Each division strictly reduces the number
    \item Total operations per number are proportional to the number of prime factors
\end{itemize}

\subsubsubsection{Time Complexity}

\[
\text{Preprocessing: } O(N \log \log N)
\]
\[
\text{Factorization: } O(\log n) \text{ per number}
\]

This approach is fast, scalable, and safe for tight constraints.

---

% =====================================================
\subsubsection{Approach 2: Without SPF (Trial Division)}
% =====================================================

\subsubsubsection{Idea}

\begin{itemize}
    \item Generate all primes up to $\sqrt{10^7}$
    \item Factorize each number by trying division with all primes
    \item Accumulate and cancel prime exponents
\end{itemize}

\subsubsubsection{Why It Is Slower}

\begin{itemize}
    \item Trial division may check many primes for each number
    \item Worst-case factorization takes $O(\sqrt{n})$
    \item Repeated work across many elements
\end{itemize}

\subsubsubsection{Time Complexity}

\[
\text{Prime generation: } O(N \log \log N)
\]
\[
\text{Factorization: } O(\sqrt{n}) \text{ per number}
\]

This approach can pass, but is close to time limits.

---

% =====================================================
\subsubsection{Comparison of Approaches}
% =====================================================

\begin{center}
\begin{tabular}{|c|c|c|}
\hline
\textbf{Aspect} & \textbf{SPF-Based} & \textbf{Trial Division} \\
\hline
Accepted & Yes & Yes \\
\hline
Factorization Time & $O(\log n)$ & $O(\sqrt{n})$ \\
\hline
Speed & Fast & Slower \\
\hline
Memory Usage & Higher & Lower \\
\hline
Ease of Implementation & Moderate & Simple \\
\hline
Recommended & \textbf{Yes} & Only for small limits \\
\hline
\end{tabular}
\end{center}

---

% =====================================================
\section{Combinatorics}
% =====================================================

% =====================================================
\subsection{Factorials}
% =====================================================

\subsubsection{Definition}
\[
n! = 1 \cdot 2 \cdot 3 \cdots n
\]

---

% =====================================================
\subsection{Precomputation of Factorials}
% =====================================================

\subsubsection{Why Precompute}
Many problems require repeated computation of combinations and permutations.
Precomputing factorials avoids recomputation and makes queries efficient.

\subsubsection{Complexity}
\[
O(N) \text{ time}, \quad O(N) \text{ space}
\]

\subsubsection{Code}
\begin{lstlisting}
vector<long long> computeFactorials(int n, long long mod) {
    vector<long long> fact(n + 1);
    fact[0] = 1;
    for (int i = 1; i <= n; i++)
        fact[i] = fact[i - 1] * i % mod;
    return fact;
}
\end{lstlisting}

---

% =====================================================
\subsection{Precomputation of Inverse Factorials}
% =====================================================

\subsubsection{Why Inverse Factorials Are Needed}

To compute combinations using:
\[
\binom{n}{r} = \frac{n!}{r!(n-r)!}
\]

we need inverses of factorial values.

A naive approach would compute:
\[
(r!)^{-1},\; ((n-r)!)^{-1}
\]
for every query.

---

\subsubsection{Why Precomputing Inverse Factorials Is Better}

\paragraph{Naive Approach}
\begin{itemize}
    \item Factorials are precomputed
    \item For every query, compute inverse of $r!$ and $(n-r)!$
    \item Each inverse costs $O(\log mod)$
\end{itemize}

\paragraph{Total Cost}
For $Q$ queries:
\[
O(Q \log mod)
\]

---

\paragraph{Optimized Approach}
\begin{itemize}
    \item Compute inverse of $n!$ once
    \item Use the relation:
    \[
    (i-1)!^{-1} = i!^{-1} \cdot i
    \]
    \item Build all inverse factorials in linear time
\end{itemize}

\paragraph{Total Cost}
\[
O(N + \log mod) \text{ preprocessing}
\]
\[
O(1) \text{ per query}
\]

This is significantly faster when there are many queries.

---

\subsubsection{Code}
\begin{lstlisting}
vector<long long> computeInverseFactorials(
    int n, long long mod, const vector<long long>& fact) {

    vector<long long> ifact(n + 1);
    ifact[n] = modpow(fact[n], mod - 2, mod);

    for (int i = n - 1; i >= 0; i--)
        ifact[i] = ifact[i + 1] * (i + 1) % mod;

    return ifact;
}
\end{lstlisting}

% -----------------------------------------------------
\subsection{Practice Problem: Christmas Tree Decoration}
% -----------------------------------------------------

\paragraph{Problem}
\textbf{D. Christmas Tree Decoration}

\paragraph{Link}
\url{https://codeforces.com/problemset/problem/2190/D}

\paragraph{Problem Summary}

There are $n$ people and $(n+1)$ boxes of decorations.
Box $0$ is a shared box, while box $i$ belongs to person $i$.

People act in a given permutation order:
each person must take one decoration either from their own box or from box $0$.
This process repeats cyclically until all decorations are used.

A permutation is called \textbf{fair} if it is possible to finish the process
without ever encountering a situation where a person needs to take a decoration
but both box $0$ and their own box are empty.

Your task is to count the number of fair permutations.

---

\subsubsection{Key Observation}

Each person $i$ can safely act as long as:
\begin{itemize}
    \item They still have decorations in their own box, or
    \item Box $0$ has decorations left
\end{itemize}

This creates a global dependency on:
\begin{itemize}
    \item The total number of decorations
    \item How many people are forced to rely on box $0$
\end{itemize}

The order in which people appear matters, but the process can be analyzed
by grouping people based on how many times they must draw from box $0$.

---

\subsubsection{Core Idea}

\begin{itemize}
    \item Think in terms of \textbf{how many times each person needs help from box $0$}
    \item The process is valid if box $0$ never runs out before all forced usages are satisfied
    \item The problem reduces to counting permutations where these constraints hold
\end{itemize}

This transforms the problem into a \textbf{counting permutations with restrictions},
which is a classic combinatorics setup.

---

\subsubsection{Approach}

\begin{enumerate}
    \item Count how many decorations must come from box $0$
    \item Determine how many people are \emph{forced} to rely on box $0$
    \item Count the number of permutations where these people appear
          in an order that never exhausts box $0$
    \item Use combinations and factorials to count valid arrangements
\end{enumerate}

All counting is done modulo $998244353$.

---


% =====================================================
\subsection{Computing $nCr$ and $nPr$ Using Precomputed Arrays}
% =====================================================

\subsubsection{$nCr$}
\[
\binom{n}{r} = \frac{n!}{r!(n-r)!}
\]

\begin{lstlisting}
long long nCr(int n, int r,
              const vector<long long>& fact,
              const vector<long long>& ifact,
              long long mod) {
    if (r < 0 || r > n) return 0;
    return fact[n] * ifact[r] % mod * ifact[n - r] % mod;
}
\end{lstlisting}

\subsubsection{Time Complexity}
\[
O(1)
\]

---

\subsubsection{$nPr$}
\[
P(n,r) = \frac{n!}{(n-r)!}
\]

\begin{lstlisting}
long long nPr(int n, int r,
              const vector<long long>& fact,
              const vector<long long>& ifact,
              long long mod) {
    if (r < 0 || r > n) return 0;
    return fact[n] * ifact[n - r] % mod;
}
\end{lstlisting}

---

% =====================================================
\subsection{Computing $nCr$ Without Precomputation}
% =====================================================

\subsubsection{Idea}
Use the multiplicative form of combinations:
\[
\binom{n}{r}
=
\frac{n(n-1)(n-2)\cdots(n-r+1)}{r!}
\]

Always compute with:
\[
r = \min(r, n-r)
\]

---

\subsubsection{Code}
\begin{lstlisting}
long long nCr(int n, int r) {
    if (r < 0 || r > n) return 0;
    r = min(r, n - r);

    long long res = 1;
    for (int i = 1; i <= r; i++) {
        res = res * (n - r + i) / i;
    }
    return res;
}
\end{lstlisting}

% -----------------------------------------------------
\subsection{Practice Problem: Unfair Game}
% -----------------------------------------------------

\paragraph{Problem}
\textbf{D. Unfair Game}

\paragraph{Link}
\url{https://codeforces.com/contest/2162/problem/D}

\paragraph{Problem Summary}

Bob chooses a number $a$ from $1$ to $n$, where $n = 2^d$.
Alice can repeatedly:
\begin{itemize}
    \item Divide the number by $2$ if it is even, or
    \item Subtract $1$
\end{itemize}

After each move, Alice learns the exponent of the largest power of $2$
dividing the current number.

Alice wants to reach $0$ within at most $k$ moves.
Your task is to count how many starting numbers Alice \textbf{cannot} win from
within $k$ moves, assuming optimal play.

---

\subsubsection{Key Observation}

For any number $a > 0$, there exists a unique value:
\[
x = v_2(a)
\]
such that:
\[
a \equiv 0 \pmod{2^x}, \quad a \not\equiv 0 \pmod{2^{x+1}}
\]

This is exactly the number of trailing zeros in the binary representation of $a$.

---

\subsubsection{Core Idea}

\begin{itemize}
    \item Each division by $2$ removes one trailing zero
    \item Subtracting $1$ turns an odd number into an even one
    \item The total number of moves needed depends only on the binary structure of $a$
\end{itemize}

Thus, the game reduces to:
\begin{itemize}
    \item Analyzing binary lengths
    \item Counting how many numbers have a given number of trailing zeros
\end{itemize}

---

\subsubsection{Why Combinations Appear}

Numbers from $1$ to $2^d$ can be classified by:
\begin{itemize}
    \item Total bit length
    \item Positions of $1$ bits
\end{itemize}

Counting how many numbers fall into each category involves:
\[
\binom{d}{k}
\]

This allows us to count valid starting numbers using combinations,
without needing factorial precomputation due to small effective ranges.

---

\subsubsection{Approach}

\begin{enumerate}
    \item Let $d = \log_2(n)$
    \item Iterate over possible binary lengths
    \item For each length, count numbers requiring more than $k$ moves
    \item Use $nCr$ to count configurations efficiently
\end{enumerate}

All calculations are done using direct multiplicative $nCr$.

---

% -----------------------------------------------------
\subsection{Practice Problem: Production of Snowmen}
% -----------------------------------------------------

\paragraph{Problem}
\textbf{C. Production of Snowmen}

\paragraph{Link}
\url{https://codeforces.com/contest/2182/problem/C}

\paragraph{Problem Summary}

You are given three cyclic arrays:
\[
A = (a_1,\dots,a_n),\;
B = (b_1,\dots,b_n),\;
C = (c_1,\dots,c_n)
\]

A snowman is \textbf{stable} if:
\[
a < b < c
\]

You must choose starting indices $(i, j, k)$ such that for every shift
$t = 0 \dots n-1$:
\[
a_{i+t} < b_{j+t} < c_{k+t}
\]

The task is to count how many such triples $(i,j,k)$ exist.

---

\subsubsection{Key Observation}

Because the conveyors are \textbf{cyclic}, once $(i,j,k)$ is chosen:
\begin{itemize}
    \item All comparisons are shift-invariant
    \item Only relative alignment matters, not absolute position
\end{itemize}

Thus, the problem reduces to checking consistency of inequalities
across all $n$ positions.

---

\subsubsection{Reformulation}

For a fixed offset alignment:
\[
(a_x, b_y, c_z)
\]

we need:
\[
\forall t:\quad
a_{x+t} < b_{y+t} \;\text{and}\; b_{y+t} < c_{z+t}
\]

This splits into two independent conditions:
\begin{itemize}
    \item Valid alignments of $A$ with $B$
    \item Valid alignments of $B$ with $C$
\end{itemize}

---

\subsubsection{Counting Strategy}

\begin{itemize}
    \item Precompute all valid cyclic shifts where $a_i < b_j$
    \item Precompute all valid cyclic shifts where $b_j < c_k$
    \item Combine compatible shifts using counting
\end{itemize}

Each valid configuration contributes multiplicatively.

---

% =====================================================
\subsection{Binomial Theorem}
% =====================================================

\subsubsection{Statement}
\[
(a + b)^n = \sum_{k=0}^{n} \binom{n}{k} a^{n-k} b^k
\]

\subsubsection{General Term}
\[
T_{k+1} = \binom{n}{k} a^{n-k} b^k
\]

\subsubsection{Finding Coefficients}
\begin{itemize}
    \item Coefficient of $a^{n-k}b^k$: $\binom{n}{k}$
    \item Coefficient of $x^k$ in $(ax+b)^n$: $\binom{n}{k}a^k b^{n-k}$
\end{itemize}

---

% =====================================================
\subsection{Stars and Bars}
% =====================================================

\subsubsection{Problem Statement}
Count the number of non-negative integer solutions to:
\[
x_1 + x_2 + \cdots + x_n = k
\]

---

\subsubsection{Intuition}

Represent $k$ identical objects as stars:
\[
\star \star \star \cdots \star
\]

To divide them into $n$ groups, place $n-1$ bars between stars.
Each distinct arrangement corresponds to a unique solution.

---

\subsubsection{Formula}
\[
\boxed{\binom{n + k - 1}{k}}
\]

---

\subsubsection{Why This Works}
\begin{itemize}
    \item Total positions: $k + n - 1$
    \item Choose $k$ positions for stars
    \item Remaining positions are bars
\end{itemize}

---

% -----------------------------------------------------
\subsection{Practice Problem: Stars \& Bars}
% -----------------------------------------------------

\paragraph{Problem}
\textbf{CF 1288C — Two Arrays}

\paragraph{Link}
\url{https://codeforces.com/problemset/problem/1288/C}

\paragraph{Core Idea}

The problem asks for the number of ways to build an array with fixed sum
constraints on all prefixes and total sum.

Key combinatorial insight:
\begin{itemize}
    \item You are effectively distributing a fixed total sum across positions
    \item Each prefix has a minimum required sum
    \item This leads to counting integer solutions with lower bounds
\end{itemize}

This is a typical **Stars and Bars with constraints** scenario, where
we adjust variables to account for minimum sums, then apply standard counting.

\paragraph{Why It Fits}

Stars and Bars solves counts of non-negative integer solutions:
\[
x_1 + x_2 + \cdots + x_n = k
\]
But when there are **prefix minimum constraints**, you:
\begin{itemize}
    \item shift variables to remove lower bounds
    \item transform into a standard root counting problem
    \item apply combination formulas accordingly
\end{itemize}

---

% =====================================================
\section{Probability and Expectation}
% =====================================================

\subsection{Probability Basics}

\subsubsection{Definition}
\[
P(A) = \frac{\text{favorable outcomes}}
             {\text{total outcomes}}
\]

In competitive programming:
\[
P \equiv \text{favorable} \times inv(\text{total}) \pmod{mod}
\]

---

\subsection{Probability Under Modulo}

Division is replaced by modular inverse:
\[
\frac{a}{b} \equiv a \cdot b^{mod-2} \pmod{mod}
\]

---

\subsection{Conditional Probability}

\[
P(A \mid B) = \frac{P(A \cap B)}{P(B)}
\]

Used when outcomes depend on previous events.

---

\subsection{Independent Events}

Events $A$ and $B$ are independent if:
\[
P(A \cap B) = P(A) \cdot P(B)
\]

---

\subsection{Expected Value}

\subsubsection{Definition}
\[
\mathbb{E}[X] = \sum x \cdot P(X = x)
\]

---

\subsection{Linearity of Expectation}

\[
\mathbb{E}[X_1 + X_2 + \cdots + X_n]
=
\sum_{i=1}^{n} \mathbb{E}[X_i]
\]

---

\subsection{Probability DP}

\[
dp[i] = \sum dp[j] \cdot P(j \rightarrow i)
\]

---

% -----------------------------------------------------
\subsection{Practice Problem (Probability)}
% -----------------------------------------------------

\paragraph{Problem}
\textbf{CF 2125D — Segments Covering}

\paragraph{Link}
\url{https://codeforces.com/problemset/problem/2125/D}

\paragraph{Goal}
Compute the probability that \textbf{each cell is covered by exactly one segment},
where each segment exists independently with probability $\frac{p_i}{q_i}$.

\subsubsection{Key Probability Transformation}

For each segment with probability
\[
P_i = \frac{p_i}{q_i},
\]
rewrite it as:
\[
P_i = (1 - P_i) \cdot \frac{P_i}{1 - P_i}.
\]

This allows the total probability to be expressed as:
\[
\left( \prod_i (1 - P_i) \right)
\cdot
\left( \sum \text{valid configurations using } \frac{P_i}{1 - P_i} \right).
\]

\begin{itemize}
    \item The product term accounts for all segments initially being inactive
    \item The summation term activates exactly one valid segment per cell
\end{itemize}

\subsubsection{DP Interpretation}

\begin{itemize}
    \item Segments are grouped by their right endpoint
    \item Each segment contributes a weight:
    \[
    w_i = \frac{P_i}{1 - P_i}
    \]
    \item A suffix DP is computed from right to left
    \item DP ensures that:
    \begin{itemize}
        \item exactly one segment covers each cell
        \item overlapping segments are handled without double counting
    \end{itemize}
\end{itemize}

The DP value at position $i$ represents the total contribution of valid
segment selections covering suffix $[i, m-1]$.

\subsubsection{Why Reverse (Suffix) DP Works}

Processing from right to left ensures:
\begin{itemize}
    \item each segment contributes only to the cells it covers
    \item overlapping probabilities are merged correctly
    \item independence between segments is preserved algebraically
\end{itemize}

This avoids explicit enumeration of cells or subsets.

\subsubsection{Final Formula}

The final probability is:
\[
\text{Answer} =
\left( \prod_i (1 - P_i) \right)
\cdot
\text{DP}[0]
\pmod{998244353}.
\]

\subsubsection{Complexity}

\begin{itemize}
    \item Time Complexity: $O(n + m)$
    \item Space Complexity: $O(n + m)$
\end{itemize}


% =====================================================
\section{Series and Summations}
% =====================================================

\subsection{Arithmetic Progression (AP)}

\subsubsection{Nth Term}
\[
a_n = a + (n-1)d
\]

\subsubsection{Sum of First $n$ Terms}
\[
S_n = \frac{n}{2} (2a + (n-1)d)
\]

\subsubsection{Special Case}
\[
1 + 2 + \cdots + n = \frac{n(n+1)}{2}
\]

% -----------------------------------------------------
\subsection{Practice Problem (Arithmetic Progression)}
% -----------------------------------------------------

\paragraph{Problem}
\textbf{CF 1616C — Representative Edges}

\paragraph{Link}
\url{https://codeforces.com/contest/1616/problem/C}

\paragraph{Key Observation}
An array is an arithmetic progression if and only if there exist constants
$a$ and $d$ such that:
\[
a_i = a + i \cdot d
\]

Equivalently, for any two indices $i < j$:
\[
d = \frac{a_j - a_i}{j - i}
\]

---

\paragraph{Approach}
\begin{itemize}
    \item Fix two indices $(i, j)$ to uniquely determine the common difference $d$
    \item This fixes the entire arithmetic progression
    \item For all positions $k$, check whether:
    \[
    a_k = a_i + (k - i)\cdot d
    \]
    \item Count how many elements already match this AP
    \item Minimum operations $=$ total elements $-$ maximum matches
\end{itemize}

---

\subsection{Geometric Progression (GP)}

\subsubsection{Nth Term}
\[
a_n = a r^{n-1}
\]

\subsubsection{Sum ($r \neq 1$)}
\[
S_n = a \cdot \frac{r^n - 1}{r - 1}
\]

\subsubsection{Infinite GP ($|r| < 1$)}
\[
S = \frac{a}{1-r}
\]

% -----------------------------------------------------
\subsection{Practice Problem (Geometric Progression)}
% -----------------------------------------------------

\paragraph{Problem}
\textbf{CF 2182B — New Year Cake}

\paragraph{Link}
\url{https://codeforces.com/contest/2182/problem/B}

\paragraph{Key Observation}
The sizes of cake layers follow a geometric progression:
\[
1,\; 2,\; 4,\; 8,\; \dots
\]
Each layer is twice the size of the layer above.

Additionally:
\begin{itemize}
    \item Chocolate types must alternate
    \item The amount of chocolate required for a layer equals its size
\end{itemize}

---

\paragraph{Approach}
\begin{itemize}
    \item Consider two cases:
    \begin{itemize}
        \item Top layer starts with white chocolate
        \item Top layer starts with dark chocolate
    \end{itemize}
    \item For a fixed number of layers $k$:
    \begin{itemize}
        \item Sum the GP terms assigned to white layers
        \item Sum the GP terms assigned to dark layers
    \end{itemize}
    \item Check whether both sums fit within available chocolate
    \item Maximize $k$
\end{itemize}

---

\subsection{Common Summation Formulas}

\[
\sum_{i=1}^n i^2 = \frac{n(n+1)(2n+1)}{6}
\]

\[
\sum_{i=1}^n i^3 = \left( \frac{n(n+1)}{2} \right)^2
\]

---

\subsection{Telescoping Series}

\subsubsection{Definition}
A series where most terms cancel out.

\subsubsection{Example}
\[
\sum_{k=1}^n \left( \frac{1}{k} - \frac{1}{k+1} \right)
= 1 - \frac{1}{n+1}
\]

---

\subsubsection{General Form}
If:
\[
a_k = f(k) - f(k+1)
\]

Then:
\[
\sum_{k=1}^n a_k = f(1) - f(n+1)
\]

---

% =====================================================
\subsection{Prefix Sums}
% =====================================================

\subsubsection{Definition}

Given an array:
\[
a_1, a_2, \dots, a_n
\]

The \textbf{prefix sum array} is defined as:
\[
p_i = a_1 + a_2 + \cdots + a_i
\]
with:
\[
p_0 = 0
\]

---

\subsubsection{Why Prefix Sums Work}

Using prefix sums, any subarray sum can be computed in constant time:
\[
\sum_{i=l}^{r} a_i = p_r - p_{l-1}
\]

This converts repeated summation from:
\[
O(n) \rightarrow O(1)
\]

after a single preprocessing step.

---

\subsubsection{Time and Space Complexity}

\begin{itemize}
    \item Preprocessing: $O(n)$
    \item Each range query: $O(1)$
    \item Extra space: $O(n)$
\end{itemize}

---

% =====================================================
\subsection{Prefix XOR}
% =====================================================

\subsubsection{Definition}

For bitwise XOR:
\[
px_i = a_1 \oplus a_2 \oplus \cdots \oplus a_i
\]

with:
\[
px_0 = 0
\]

---

\subsubsection{Key Property}

Using XOR:
\[
a_l \oplus a_{l+1} \oplus \cdots \oplus a_r
=
px_r \oplus px_{l-1}
\]

This works because:
\[
x \oplus x = 0
\]

---

% -----------------------------------------------------
\subsubsection{Practice Problem (Prefix XOR Construction)}
% -----------------------------------------------------

\paragraph{Problem}
\textbf{CF 2175B — XOR Array}

\paragraph{Link}
\url{https://codeforces.com/contest/2175/problem/B}

\paragraph{Problem Insight}

The XOR of a subarray can be expressed using prefix XOR:
\[
a_l \oplus a_{l+1} \oplus \cdots \oplus a_r
=
px_r \oplus px_{l-1}
\]

Thus, enforcing XOR conditions on subarrays translates into constraints
on prefix XOR values.

---

\paragraph{Core Idea}

The problem requires:
\begin{itemize}
    \item Exactly one subarray $(l, r)$ whose XOR is $0$
    \item All other subarrays must have non-zero XOR
\end{itemize}

Using prefix XOR:
\[
px_r = px_{l-1}
\]
must hold for exactly one pair of indices, and all other prefix XOR pairs
must be distinct.

---

\paragraph{Construction Strategy}

\begin{itemize}
    \item Assign prefix XORs such that:
    \begin{itemize}
        \item $px_{l-1} = px_r$
        \item All other prefix XOR values are unique
    \end{itemize}
    \item Build the array using:
    \[
    a_i = px_i \oplus px_{i-1}
    \]
\end{itemize}

This guarantees:
\begin{itemize}
    \item XOR of subarray $(l, r)$ is $0$
    \item No other subarray XOR becomes $0$
\end{itemize}

---


\subsection{Suffix Sums and XOR}
% =====================================================

\subsubsection{Definition}

\textbf{Suffix Sum} is defined as:
\[
s_i = a_i + a_{i+1} + \cdots + a_n
\]

\textbf{Suffix XOR} is defined as:
\[
sx_i = a_i \oplus a_{i+1} \oplus \cdots \oplus a_n
\]

Suffix sums and suffix XOR work in the same way as prefix sums and prefix XOR,
but are computed from the end of the array instead of the beginning.

\subsubsection{Relation with Prefix Arrays}

If $p_i$ is the prefix sum and $px_i$ is the prefix XOR, then:
\[
s_i = p_n - p_{i-1}
\]

\[
sx_i = px_n \oplus px_{i-1}
\]

---

% =====================================================
\section{Pigeonhole Principle}
% =====================================================

\subsection{Statement}

\subsubsection{Basic Form}
If more than $n$ objects are placed into $n$ boxes, then at least one box
contains more than one object.

Formally:
\[
\text{If } k > n \text{ objects are distributed into } n \text{ boxes,}
\Rightarrow \exists \text{ a box with at least } 2 \text{ objects.}
\]

---

\subsubsection{Generalized Form}
If $k$ objects are placed into $n$ boxes, then at least one box contains at least:
\[
\left\lceil \frac{k}{n} \right\rceil
\]
objects.

---

\subsection{Why This Works (Intuition)}

Assume the opposite:
\begin{itemize}
    \item Every box contains at most $\left\lfloor \frac{k}{n} \right\rfloor$ objects
    \item Then total objects $\le n \cdot \left\lfloor \frac{k}{n} \right\rfloor < k$
\end{itemize}

This contradicts the assumption that $k$ objects exist.

Hence, the principle must hold.

---

\subsection{Simple Examples}

\subsubsection{Example 1}
Among 13 people, at least two people have birthdays in the same month.

\paragraph{Explanation}
\begin{itemize}
    \item Objects: 13 people
    \item Boxes: 12 months
    \item Since $13 > 12$, some month has at least 2 birthdays
\end{itemize}

---

\subsubsection{Example 2}
If 11 integers are chosen from $\{1,2,\dots,10\}$, at least one number is repeated.

---

\subsection{Classic CP Use Case}

\paragraph{Subarray Sum Divisible by $n$}

Given an array of length $n$, consider prefix sums modulo $n$:
\[
p_i = (a_1 + a_2 + \cdots + a_i) \bmod n
\]

There are $n+1$ values:
\[
p_0, p_1, \dots, p_n
\]

But only $n$ possible remainders.

By the Pigeonhole Principle:
\begin{itemize}
    \item Either some $p_i = 0$
    \item Or two prefix sums are equal
\end{itemize}

In both cases, a subarray with sum divisible by $n$ exists.

---
% -----------------------------------------------------
\subsection{Practice Problem (Pigeonhole Principle)}
% -----------------------------------------------------

\paragraph{Problem}
\textbf{CF 2120D — Matrix Game}

\paragraph{Link}
\url{https://codeforces.com/contest/2120/problem/D}

\paragraph{Main Idea}

Each cell of the matrix can take one of $k$ values.
Harshith tries to avoid creating an $a \times b$ submatrix with all equal elements.

\begin{itemize}
    \item Consider all possible placements of values in the matrix
    \item The number of submatrices of size $a \times b$ grows with $n$ and $m$
    \item The number of distinct value patterns Harshith can enforce is limited by $k$
\end{itemize}

When the matrix becomes large enough, the number of required configurations
exceeds the number of distinct value choices.

By the \textbf{Pigeonhole Principle}, some $a \times b$ submatrix must have
all equal elements, guaranteeing Aryan’s win.

\paragraph{Approach}

\begin{itemize}
    \item Observe that Harshith controls the matrix values, but each cell can take
    only one of $k$ distinct values
    \item Aryan wins if there exists an $a \times b$ submatrix with all equal values
    \item Count the number of distinct $a \times b$ submatrices in an $n \times m$ matrix
    \item Compare this with the maximum number of distinct value assignments Harshith
    can use to avoid repetition
\end{itemize}

Once the number of required submatrices exceeds the number of distinct patterns
possible using $k$ values, avoiding repetition becomes impossible.

By the \textbf{Pigeonhole Principle}, at least one $a \times b$ submatrix
must have all elements equal, and Aryan is guaranteed to win.


\paragraph{Key Insight}

This is an extremal construction problem:
\begin{itemize}
    \item If objects (submatrices) $>$ boxes (distinct value patterns)
    \item Then repetition is unavoidable
\end{itemize}

The task is to find the smallest $(n,m)$ where this inevitability occurs.

% =====================================================
\section{Basic Geometry for Competitive Programming}
% =====================================================

\subsection{Distance Between Two Points}

Let two points be:
\[
A(x_1, y_1), \quad B(x_2, y_2)
\]

---

\subsubsection{Euclidean Distance}

\paragraph{Formula}
\[
d = \sqrt{(x_1 - x_2)^2 + (y_1 - y_2)^2}
\]

\paragraph{Important Trick}
Instead of comparing distances, compare \textbf{squared distances}:
\[
(x_1 - x_2)^2 + (y_1 - y_2)^2
\]
to avoid floating-point inaccuracies.

---

\subsubsection{Manhattan Distance}

\paragraph{Formula}
\[
d = |x_1 - x_2| + |y_1 - y_2|
\]

\paragraph{Interpretation}
Distance when movement is restricted to horizontal and vertical steps.

% -----------------------------------------------------
\subsection{Practice Problem (Geometry)}
% -----------------------------------------------------

\paragraph{Problem}
\textbf{CF 2119B — Line Segments}

\paragraph{Link}
\url{https://codeforces.com/contest/2119/problem/B}

\paragraph{Main Idea}

Each move allows you to travel \emph{exactly} a fixed distance, but in any direction.
Thus, after all operations, the question reduces to whether the final point can be
reached using segments of given lengths.

Let:
\begin{itemize}
    \item $D$ be the Euclidean distance between the start and end points
    \item $S = a_1 + a_2 + \cdots + a_n$ be the total available movement
    \item $M = \max(a_1, a_2, \ldots, a_n)$ be the largest single move
\end{itemize}

\paragraph{Approach}

Using the triangle inequality:
\begin{itemize}
    \item The farthest reachable distance is $S$
    \item The closest reachable distance is $\max(0,\, 2M - S)$
\end{itemize}

Reaching the target point is possible if and only if:
\[
\max(0,\, 2M - S) \le D \le S
\]

\paragraph{Why It Works}

Each move can be rotated arbitrarily, so only the total length and imbalance between
segments matters. This is a direct application of geometric feasibility using
triangle inequality constraints.
---

% =====================================================
\section{Basic Statistics}
% =====================================================

\subsection{Mean}

\subsubsection{Definition}
\[
\text{Mean} = \frac{x_1 + x_2 + \cdots + x_n}{n}
\]

\subsubsection{Uses}
\begin{itemize}
    \item Expected value problems
    \item Average operations
    \item Probability calculations
\end{itemize}

---

\subsection{Median}

\subsubsection{Definition}
The median is the middle element after sorting the array.

\begin{itemize}
    \item If $n$ is odd: element at position $\frac{n}{2}$
    \item If $n$ is even: any value between the two middle elements
\end{itemize}

---

\subsubsection{Key Property}
The median minimizes:
\[
\sum |x_i - m|
\]

This makes it extremely important in optimization problems.
---

% -----------------------------------------------------
\subsection{Practice Problem (Median)}
% -----------------------------------------------------

\paragraph{Problem}
\textbf{CF 2149D — A and B}

\paragraph{Link}
\url{https://codeforces.com/contest/2149/problem/D}

\paragraph{Main Idea}

We want all characters of one type (either all \texttt{a} or all \texttt{b})
to form a single contiguous block using the minimum number of adjacent swaps.

This is equivalent to:
\begin{itemize}
    \item Fixing a target block of consecutive positions
    \item Moving all chosen characters into this block
    \item Minimizing the total number of swaps (movement distance)
\end{itemize}

\paragraph{Key Observation}

Adjacent swaps correspond to \textbf{absolute movement cost}.
Thus, minimizing swaps is the same as minimizing:
\[
\sum |pos_i - target_i|
\]

This sum is minimized when the block is centered at the \textbf{median}
of the chosen character positions.

\paragraph{Approach}

\begin{itemize}
    \item Collect indices of all \texttt{a}'s (or all \texttt{b}'s)
    \item Assume they form a contiguous block
    \item Align them to consecutive positions
    \item The optimal alignment is achieved by placing the block around the median
    \item Compute total movement cost
    \item Do this for both characters and take the minimum
\end{itemize}

\paragraph{Why It Works}

The median minimizes the sum of absolute differences.
Since each adjacent swap shifts a character by one position,
the total number of swaps is minimized by median alignment.

% -----------------------------------------------------
\subsection{Practice Problem (Median \& Manhattan Distance)}
% -----------------------------------------------------

\paragraph{Problem}
\textbf{CF 2122C — Manhattan Pairs}

\paragraph{Link}
\url{https://codeforces.com/problemset/problem/2122/C}

\paragraph{Main Idea}

We are given $n$ points on a 2D plane ($n$ is even), and must form $\frac{n}{2}$
disjoint pairs to \textbf{maximize the total Manhattan distance}:
\[
\sum |x_a - x_b| + |y_a - y_b|
\]

The key observation is that Manhattan distance is separable across dimensions,
and optimal pairing is achieved by pairing \textbf{far-apart points}.

\paragraph{Key Observations}

\begin{itemize}
    \item Sorting points by $x$ separates them spatially
    \item Pairing smallest $x$ with largest $x$ maximizes $|x_a - x_b|$
    \item After fixing $x$ separation, pairing extremes in $y$ further
          maximizes $|y_a - y_b|$
    \item This mirrors the idea that distance is maximized when elements are
          placed on opposite sides of the median
\end{itemize}

\paragraph{Greedy Strategy}

\begin{itemize}
    \item Sort all points by $x$
    \item Split into two halves:
    \begin{itemize}
        \item Left half = smaller $x$
        \item Right half = larger $x$
    \end{itemize}
    \item Sort both halves by $y$
    \item Pair:
    \[
        \text{left}[i] \;\; \text{with} \;\; \text{right}[\tfrac{n}{2}-1-i]
    \]
    \item This pairs smallest $y$ with largest $y$, maximizing vertical distance
\end{itemize}

\paragraph{Why This Works}

For Manhattan distance:
\[
|x_1 - x_2| + |y_1 - y_2|
\]
both components are maximized independently by pairing elements symmetrically
around their medians.

---

\subsection{Mode}

\subsubsection{Definition}
The mode is the most frequent element in a dataset.

\subsubsection{Uses}
\begin{itemize}
    \item Frequency-based decisions
    \item Majority element problems
    \item Hash map counting problems
\end{itemize}

% -----------------------------------------------------
\subsection{Practice Problem: Majority Element}
% -----------------------------------------------------

\paragraph{Problem}
\textbf{LeetCode — Majority Element}

\paragraph{Link}
\url{https://leetcode.com/problems/majority-element/description/}

\paragraph{Core Idea}

A “majority element” is an element that appears **more than $\lfloor n/2 \rfloor$ times** in an array of length $n$.

Key observations:
\begin{itemize}
    \item The majority element is unique if it exists.
    \item A naive frequency map takes $O(n)$ time and $O(n)$ space.
    \item There exists an $O(n)$ time and $O(1)$ space algorithm:
          the \textbf{Boyer–Moore Voting Algorithm}.
\end{itemize}

Boyer–Moore works by maintaining a candidate and a counter:
\[
\text{if counter } = 0, \text{ set candidate = current element}
\]
\[
\text{if current == candidate, counter++ else counter--}
\]

The final candidate is the majority element if a majority element exists (since it survives net cancellation).


% -----------------------------------------------------
\subsection{Practice Problem: Majority Element II}
% -----------------------------------------------------

\paragraph{Problem}
\textbf{LeetCode — Majority Element II}

\paragraph{Link}
\url{https://leetcode.com/problems/majority-element-ii/description/}

\paragraph{Core Idea}

This is a generalization: find all elements that appear **more than $\lfloor n/3 \rfloor$ times**.

There can be at most two such elements.
The optimized $O(n)$, $O(1)$ space method extends Boyer–Moore:

\begin{itemize}
    \item Maintain two candidates and two counters.
    \item First pass: adjust candidates/counters by cancellation.
    \item Second pass: verify actual counts.
\end{itemize}

---

% -----------------------------------------------------
\section{Additional Practice Problems}
% -----------------------------------------------------

\paragraph{Suffix Sums \& Combinatorics :}
\textbf{CF 2008F} \\
\url{https://codeforces.com/contest/2008/problem/F}

\paragraph{Bitwise Operations :}
\textbf{CF 1614C} \\
\url{https://codeforces.com/problemset/problem/1614/C}

\paragraph{Factors :}
\textbf{CF 1925B} \\
\url{https://codeforces.com/contest/1925/problem/B}

\paragraph{Combinatorics :}
\textbf{CF 2150B} \\
\url{https://codeforces.com/problemset/problem/2150/B}

\paragraph{Brute Force :}
\textbf{CF 2144D} \\
\url{https://codeforces.com/problemset/problem/2144/D}

\paragraph{Maths (Log2) :}
\textbf{CF 2013B} \\
\url{https://codeforces.com/problemset/problem/2109/B}

\paragraph{Binary Numbers :}
\textbf{CF 1225C} \\
\url{https://codeforces.com/problemset/problem/1225/C}

\paragraph{Factors :}
\textbf{CF 2193E} \\
\url{https://codeforces.com/problemset/problem/2193/E}

\paragraph{Bitwise :}
\textbf{CF 2189C2} \\
\url{https://codeforces.com/problemset/problem/2189/C2}

\paragraph{Telescopic Series :}
\textbf{CF 2183E} \\
\url{https://codeforces.com/problemset/problem/2183/E}

\paragraph{GCD : }
\textbf{CF 2116C} \\
\url{https://codeforces.com/contest/2116/problem/C}

\paragraph{Digits :}
\textbf{CF 2132D} \\
\url{https://codeforces.com/problemset/problem/2132/D}

\paragraph{Prefix/ Suffix GCD :}
\textbf{CF 2126E} \\
\url{https://codeforces.com/problemset/problem/2126/E}

\paragraph{Arithemetic Progression :}
\textbf{CF 2117D} \\
\url{https://codeforces.com/problemset/problem/2117/D}

\paragraph{Number Theory}
\textbf{CF 2109C2} \\
\url{https://codeforces.com/problemset/problem/2109/C2}


\end{document}
